% Section 11.3, from lectures/L11/appendix/b_exercises.md.

\section{Exercises}
\label{c11:sec:exercises}\label{c11:app:b}

\begin{aside}
Nine, ending with the Cross-check. Do them in order; Exercise~\ref{c11:ex:6} onwards need the target
running.

Where an exercise can be checked mechanically, a \textbf{Check yourself} line says how.

\textbf{Solutions.} The exercises that are not Code are answered in \secref{sol:sec:c11}. The Code
exercises have no solution: each is checked by running it, as its Check yourself line says.
\end{aside}

% ------------------------------------------------------------------------------------------
\exercise[fillaccent2]{What a framework gives you}{Recall}
\label{c11:ex:1}

\xpart{a} Your Chapter~\ref{c9:ch} driver works and presents \file{/dev/qa_wait}. List five things
it does not have that a subsystem would have provided.

\xpart{b} State the trade a framework offers, in one sentence.

\xpart{c} Name three separate costs of adding an \code{ioctl} of your own. One of them only appears
on a particular kind of system; which, and what kind?

\xpart{d} \code{ls /sys/bus/iio/devices} finds every IIO device on a machine. What is the equivalent
command for finding every device that implements your own ioctl?

% ------------------------------------------------------------------------------------------
\exercise[fillaccent2]{The tour}{Recall}
\label{c11:ex:2}

For each device, name the subsystem or subsystems it belongs to.

\xpart{a} A temperature sensor on an I2C bus.

\xpart{b} A pin that drives an LED.

\xpart{c} A hardware timer that resets the board if not kicked.

\xpart{d} A rotary encoder a person turns.

\xpart{e} A NAND flash chip.

\xpart{f} A 16-bit ADC on SPI.

\xpart{g} A crystal oscillator feeding three peripherals.

Then:

\xpart{h} Two of your answers are ``more than one subsystem''. Which, and why is that not a
contradiction?

\xpart{i} Which of the ten subsystems in \S\ref{c11:app:a3} presents \textbf{nothing} to userspace,
and what is it for instead?

% ------------------------------------------------------------------------------------------
\exercise[fillaccent3]{Reading an IIO device}{Hand calculation}
\label{c11:ex:3}

A device exposes \code{in_voltage0_raw} and \code{in_voltage0_scale}. Reading them gives:

\begin{console}
in_voltage0_raw    = 2048
in_voltage0_scale  = 0.805664
\end{console}

\xpart{a} What is the voltage? What are the units of \code{scale}, and how do you know?

\xpart{b} Why does IIO split the value into a raw count and a scale, rather than having the driver
report volts directly?

\xpart{c} The driver adds \code{IIO_CHAN_INFO_OFFSET} and it reads \code{-100}. Recompute the
voltage. In what order are raw, offset and scale applied?

\xpart{d} Your device reports a plain counter with no physical meaning at all. Which of \code{raw},
\code{scale} and \code{offset} should the driver implement, and what should it do about \code{type}?

% ------------------------------------------------------------------------------------------
\exercise[fillaccent4]{The pin that is not yet a GPIO}{Design}
\label{c11:ex:4}

\xpart{a} Before a pin can be driven as a GPIO, something must have multiplexed it as one rather
than as a UART line. Which subsystem, and where is the configuration written?

\xpart{b} Should your driver perform that multiplexing in \code{probe}? Give the reason.

\xpart{c} Your device needs a clock enabled before its registers respond. Write the two calls, and
say what happens if you instead write the SoC's clock gate register directly and a second driver
shares that clock.

\xpart{d} Your device has a regulator that must be on. Same question: what does the framework do
that a direct register write does not?

\xpart{e} All three of these are dependencies that may not be ready when your \code{probe} runs.
What should \code{probe} return, and what does the core do with it?

% ------------------------------------------------------------------------------------------
\exercise[fillaccent4]{When a character device is right}{Design}
\label{c11:ex:5}

\xpart{a} Give a device that genuinely fits no subsystem, and say why.

\xpart{b} For it, would you register the long way as in Chapter~\ref{c5:ch}, or use
\code{misc_register}? Justify it.

\xpart{c} You add two \code{ioctl}s. Write down what you now owe the people who use them, and for
how long.

\xpart{d} A colleague proposes an ioctl that returns a \code{struct} containing a \code{long} and a
pointer. Name both problems.

\xpart{e} State the difference between ``no subsystem fits'' and ``implementing a subsystem
interface looked like more work''. Which is a reason and which is an excuse?

% ------------------------------------------------------------------------------------------
\exercise{Register with two frameworks}{Code}
\label{c11:ex:6}

Rewrite your Chapter~\ref{c10:ch} driver as \file{lectures/L11/lab/qa_frameworks.c}.

\textbf{IIO:}
\begin{itemize}
  \item Allocate with \code{devm_iio_device_alloc}, sized for your private structure, and get it
    back with \code{iio_priv}.
  \item One channel: \code{IIO_VOLTAGE}, indexed, channel 0, with \code{IIO_CHAN_INFO_RAW}.
  \item \code{read_raw} returns the most recent sample your interrupt handler stored, as
    \code{IIO_VAL_INT}.
  \item Name the device \code{qa_dev}, mode \code{INDIO_DIRECT_MODE}, and register with
    \code{devm_iio_device_register}.
\end{itemize}
\textbf{GPIO:}
\begin{itemize}
  \item A \code{gpio_chip} with eight lines, labelled with \code{dev_name(&pdev->dev)}, \code{.base
    = -1}, parented to your platform device.
  \item \code{get}, \code{set}, \code{direction_input}, \code{direction_output} and
    \code{get_direction}. Lines 0 to 3 are outputs; lines 4 to 7 are inputs wired back to them by
    the hardware.
  \item Register with \code{devm_gpiochip_add_data}.
\end{itemize}
\textbf{And no character device}: no \code{file_operations}, no \code{read}, no \code{poll} and no
\file{/dev} node of your own, where your Chapter~\ref{c9:ch} driver had all four.

\checkyourself \code{make test L=L11} reports \textbf{PASSED}, with twenty-two checks, of which
eight come from a GPIO test program that contains no reference to your driver.

\xpart{a} \code{devm_iio_device_alloc} allocates your private structure for you. Which call from
your Chapter~\ref{c10:ch} driver does that replace, and why does the framework want to own the
allocation?

\xpart{b} Your module fails to load with \code{Unknown symbol devm_iio_device_alloc}. What is wrong, and
what is the one command that fixes it? Why is IIO a module here at all?

\xpart{c} \code{.base = -1}. What would happen on this target if you hardcoded 0 instead?

% ------------------------------------------------------------------------------------------
\exercise{Drive it with tools you did not write}{Code}
\label{c11:ex:7}

\xpart{a} Read \file{/sys/bus/iio/devices/iio:device0/name} and \code{in_voltage0_raw}. Which of
these did you write code to create?

\xpart{b} Find your gpiochip under \file{/sys/bus/gpio/devices/}. What is its parent, and what does
its \code{of_node} symlink point at? Trace that path back to Chapter~\ref{c10:ch}.

\xpart{c} Run the shipped \code{qa_gpio_test}. Read its source. Identify every line that would have
to change if it were pointed at a completely different GPIO driver on different hardware.

\xpart{d} \code{ls /sys/class/gpio} fails. Explain why that is correct rather than a missing
feature, and name the config option and the chapter that decided it.

\xpart{e} \code{libgpiod}'s \code{gpioinfo}, \code{gpioget} and \code{gpioset} are not on this
target. Would they work if they were? What exactly would they be talking to?

% ------------------------------------------------------------------------------------------
\exercise[fillaccent4]{What you gave up}{Design}
\label{c11:ex:8}

\xpart{a} Your Chapter~\ref{c9:ch} driver could block a reader until data arrived. Can a program
using \code{in_voltage0_raw} do that? What has been lost?

\xpart{b} IIO has an answer to \textbf{a)}, which this lab does not implement. Find it in the kernel
documentation and describe in two sentences what it provides and what a driver must add.

\xpart{c} Name one thing about your driver that is now \emph{harder} than it was with a character
device.

\xpart{d} Name a device for which forcing it into IIO would produce worse code than a character
device, and say what makes it a bad fit.

\xpart{e} The GPIO character device is on its second version, \code{GPIO_V2_*}. What does that tell
you about the claim that a framework removes the ABI problem? Restate the claim accurately.

% ------------------------------------------------------------------------------------------
\exercise[fillaccent]{A sample rate measured through an interface that knows nothing about it}{Cross-check}
\label{c11:ex:9}

The device is in counter mode, so sample \emph{n} has the value \emph{n}. That makes the value
readable through IIO a running total, and the difference between two reads a measurement of
something the IIO interface has no concept of.

\xpart{a} \textbf{Predict.} Your device is programmed with \code{PERIOD_NS=1000000}. Predict:
\begin{itemize}
  \item the value of \code{in_voltage0_raw} one second after the driver loads;
  \item the difference between two reads taken one second apart;
  \item what both would be at \code{period_ns=500000}.
\end{itemize}

\xpart{b} \textbf{Measure.} Read \code{in_voltage0_raw} twice, a second apart, and record both
values and the difference. Do this three times.

\xpart{c} \textbf{Reconcile.} Answer each:
\begin{itemize}
  \item The difference is not 1000. What is it, roughly, and where have you seen that number before?
    Give the chapter and the effect.
  \item Is the shortfall here the same as the one you measured in Chapter~\ref{c8:ch}, or has the
    framework added something of its own? How would you tell?
  \item You have measured a sample rate using an interface that exposes no rate, no timestamp and no
    counter. What made that possible, and what would you have had to do instead if the device were
    in \code{QA_DEV_MODE_NOISE}?
\end{itemize}

\xpart{d} \textbf{The reads are not free either.} Time a thousand reads of \code{in_voltage0_raw}
from a shell loop and from a C program. Which is faster, by how much, and what does that tell you
about the sysfs path as a way of moving samples?

\xpart{e} From \textbf{d)}, state the sample rate above which reading \code{in_voltage0_raw} in a
loop stops being a sensible way to get data out. What does IIO offer instead, and why does it exist?

\xpart{f} A colleague reports ``the IIO driver only delivers 880 samples a second when the device is
configured for 1000''. Rewrite that as an accurate sentence, naming which component each part of the
shortfall belongs to and which parts you have actually measured rather than inferred.
